% Transcription — fonds Alexandre Grothendieck, shelfmark 43, batch 3 (pages 41-60).
% Established from the facsimile archives/batches/43/batch-03.pdf (PDF page k =
% archivists' page 40+k, no cover sheet), cut from a copy of 43.pdf fetched
% through the project's own relay
% (https://grothendieck-relay.onrender.com/source/43.pdf); tiles in
% archives/tiles/43/p41-p60. Per the skill transcribe-grothendieck. The folder
% is 207 pages in eleven batches, transcribed in parallel. Batches 1, 2 and
% 9-11 were read for the folder's conventions; pp. 41-42 continue batch 2's
% last run (Galois descent for the idele class group J_C, pp. 39-40) and
% follow its notation (plain J for his double-stroked J).
%
% Pass: Opus 5.5 (claude-opus-5-5), 2026-09-26 — first pass, unchecked against the pages by a human.
%
% Dating: the folder has its own inventory date, carried verbatim in \dating{}
% with no group sentence (the skill adds the group « Jacobiennes » (43-44)
% only for undated folders). No page of this batch carries a date.
%
% Copies. No page of this batch is evidently a photocopy: all are ink on
% white « VELIN EXTRA » paper (the watermark letters show through), blue,
% blue-black and black, with orange crayon on p. 41, red crayon arrows on
% p. 57 and pencil grids on p. 58, with the pressure and colour variation of
% originals.
%
% Pages skipped: 43, a blank cream sheet. No cover in this batch.
%
% Pages summarised: none. Pages 47, 48, 57 and 58 are sheets of scattered
% formulas and are given block by block, their layout said in \note{}s;
% the diagrams of pp. 42, 44, 48, 51, 52, 55, 58, 60 are redrawn. The fast
% prose of pp. 52 (end), 54, 55 and 59 carries a heavy apparatus: the
% formulas come through, the connecting sentences often do not.
%
% His pagination: none visible on these twenty leaves. Numbered items of
% his: (i), (ii) on p. 52 (conditions on the ordered set I: saturated chains
% joining x <= y, equal lengths) and a), b) on p. 60 (the same « condition des
% chaînes », rewritten); « 1°) » at the foot of p. 57, which breaks off. No
% internal page reference. Pp. 59 and 60 appear to be in reverse order: p. 60
% ends on « et satisfont les » after announcing pairings A_c' x A_c -> A_c'c,
% and p. 59 opens mid-sentence on « une multiplication »; said in \note{}s.
%
% What the batch is about. Pp. 41-42: end of the Galois-descent run for the
% idele classes of a curve: lim H*(pi, J_C',alpha(k-bar)) = J_C'(k) by
% Mittag-Leffler, the spectral sequence collapsing to C*(Gamma, J_C(k)),
% H^1(K,G) = Ext^1(J_K, G), Ext^1(J^0_K, G), H^2(k,G) = 0, Hom(J_K, Gamma),
% « Ext^1(J^0, Gamma) = Hom » (boxed), the pairing J(k) x Ext^1(J,G) ->
% H^1(k,G) = G(k), and the map of exact sequences Ext^1(-,G) ->
% Hom(-(k), H^1(G)) for 0 -> Z -> J -> J^0 -> 0. Pp. 44-48, under his boxed
% heading « Adèles pour plusieurs dimensions » (p. 44; batch 2's cover p. 35
% reads « Adèles des préschémas »): adeles of a scheme X of higher dimension
% built from chains of points x_0 < x_1 < x_2: rings O_[x_0 x_1 ... ],
% restricted products « local », localising maps, completions O-hat_x
% tensor K, the recursive definition A^n = prod_a lim_k A^(n-1) tensor
% O_a/m_a^k with an inclusion into prod A_[x_0...x_n] « fidèlement plate »
% (p. 46), local computations in k[[s,t]] (p. 47), and a diagram of unit
% groups O-hat*_D, K-hat*_D, K* and G(K) for a surface (p. 48). Pp. 49-50:
% Matlis-type duality over a regular ring of finite dimension: D(M) =
% Hom(M, K_*) = coprod_p Hom(M, K_p), DD(M) = prod_p prod^local_(q ⊃ p)
% Hom(D_p(M), K_q), and Hom(K_p, K_q) = A_q,p flat. Pp. 51-53: the rings
% A_[p], A_[p,q] = lim A_[p] tensor A_q/q^n A_q (« on peut y remplacer q^n
% par q^(n) »), a lattice of A_q,p for chains of primes (p. 52), then an
% ordered set I with the chain condition and A_x,y in a category C. Pp.
% 53-55: X ordered set, S_X its simplicial set of chains, the A_sigma form a
% sheaf on S_X with maps phi(sigma,u); degeneracies to be neglected (keep
% strict chains); K* a cosimplicial object in an abelian category, the
% relations partial_i^n = 0 for 1 <= i <= n and the subcomplex K'^n; a
% lattice of A_[x_i ...] for x_0 < x_1 < x_2. Pp. 56-57: the complex
% A_[x_1...x_n] -> prod A_[x_0,...,x_n] -> ..., d'^2 = d''^2 = 0, d'd'' =
% d''d', the lattice for x_0 < ... < x_3, and G(K)/G(A). P. 58: a pencil
% grid for a two-dimensional local ring (A-hat, A_r-hat, K_r-hat, K and their
% unit groups) and Ext*(M, A). Pp. 60, 59: I with the « condition des
% chaînes », rings A_c for saturated chains c, maps u_1, u_2 commuting, the
% inductive system u_c,c', products A_c' x A_c -> A_c'c making the sum a
% graded ring A* with A^0 commutative, Theta in A^1 with components
% 1_[x_0,x_1], Theta^2 = sum of 1_[x_0,x_1,x_2], and the quotient by the
% two-sided ideal generated by Theta^2.
%
% Notation fixed once for the batch: J plain for his double-stroked J (pp.
% 41-42), as batch 2 does; \mathcal{O} for his O, \underline{\mathcal{O}}
% where he underlines it (non-completed local rings, and the hatted
% underlined O-hat), \hat for his hats; \mathcal{A} for his cursive A (the
% adele rings A^n, A_[...]) and A roman where he prints it; \mathfrak{p},
% \mathfrak{q}, \mathfrak{m} for his gothic letters; \mathfrak{r} on p. 58 for
% a gothic letter not read with certainty; \coprod for his upside-down
% product (⊥⊥); K_* for his starred K (the residual complex, pp. 49-52);
% \Theta as written (pp. 59-60).
%
% Readings to check first: the gloss under Ext^1 on p. 42 and « (corps
% fini) »; the indices of the O-hats on p. 44; the whole of p. 47 right
% column and its gloss « localisé en les polynômes ... »; the title word
% of p. 49; the chains of primes on p. 52 (drawn as a lozenge, given here as
% three chains); the whole prose of pp. 54, 55 and 59; « fenêtre » on p. 59;
% the glosses [i.e. (1_{x_n} Theta)_0], [i.e. d(Theta 1_{x_0})] on p. 60.
% Apparatus: about 195 \ill{} and 110 \uncertain{} in the body.

\documentclass[11pt,a4paper]{article}
% Le préambule commun à toutes les transcriptions du fonds Grothendieck.
%
% Il tient en une idée : **l'incertitude se compose, elle ne se résout pas.**
% Une transcription qui lisse un mot douteux a détruit la seule chose pour
% laquelle on la faisait — un lecteur doit pouvoir savoir, à chaque endroit,
% ce qui était sur le feuillet et ce qui a été ajouté. Les six macros qui
% suivent sont donc l'appareil critique, et non de la décoration.
%
% Les mêmes commandes sont reconnues par `scripts/render.mjs`, qui produit la
% vue de lecture du site. Une macro ajoutée ici doit l'être aussi là-bas,
% faute de quoi le rendu échouera bruyamment — ce qui est voulu : un rendu
% silencieusement faux est pire qu'un rendu absent.

\ProvidesPackage{grothendieck}[2026/08/08 Transcriptions du fonds Grothendieck]

\usepackage{amsmath,amssymb,amsthm}
\usepackage{mathrsfs}
\usepackage{stmaryrd}
% Les diagrammes commutatifs sont partout dans ce fonds : c'est la langue
% même de la géométrie algébrique de Grothendieck, et une transcription qui
% les rendrait en prose ne transcrirait rien.
\usepackage{tikz-cd}
% Français par défaut — c'est la langue des feuillets — mais l'anglais est
% chargé aussi : la traduction et le résumé sont des documents anglais, et la
% typographie française (espaces avant les deux-points) y serait une faute.
% Ces documents basculent par \selectlanguage{english} en tête de corps.
\usepackage[english,french]{babel}
\usepackage{fontspec}
\usepackage{geometry}
\geometry{a4paper,margin=25mm,marginparwidth=28mm}
\usepackage{marginnote}
\usepackage{xcolor}
% ulem, not soul. `soul`'s \st and \ul are built on a character-by-character
% scan that XeLaTeX's font machinery breaks: the document fails with an
% undefined control sequence rather than degrading. ulem does the same two jobs
% and is Unicode-engine safe. `normalem` stops it hijacking \emph, which the
% transcriptions use for Grothendieck's own emphasis.
\usepackage[normalem]{ulem}
\usepackage{hyperref}
\hypersetup{colorlinks=true,linkcolor=black,urlcolor=black,pdfborder={0 0 0}}

% --- Métadonnées du lot -----------------------------------------------------
% Reprises telles quelles par le script de rendu, qui les affiche en tête de
% la vue de lecture. Elles fixent ce que le fichier prétend couvrir : sans
% elles, une transcription flotte sans point d'attache dans un fonds de seize
% mille feuillets.

\newcommand{\folder}[1]{\gdef\grothFolder{#1}}
\newcommand{\batch}[1]{\gdef\grothBatch{#1}}
\newcommand{\leaves}[2]{\gdef\grothFirst{#1}\gdef\grothLast{#2}}
\newcommand{\foldertitle}[1]{\gdef\grothFoldertitle{#1}}
\gdef\grothFolder{?}\gdef\grothBatch{?}\gdef\grothFirst{?}\gdef\grothLast{?}\gdef\grothFoldertitle{}

% --- Le résumé --------------------------------------------------------------
%
% Il ouvre la lecture modernisée, sous le titre « Résumé », et il est composé
% en petit. Ce n'est pas un ornement : c'est le seul passage du document qui
% ne soit pas des mathématiques, et un lecteur doit pouvoir le reconnaître
% avant de l'avoir lu — puis le sauter s'il connaît déjà le sujet, ou s'y
% arrêter s'il ne le connaît pas. Le filet qui le suit marque la reprise du
% corps.

\newenvironment{resume}
  {\small\setlength{\parskip}{0.45em}}
  {\par\medskip\hrule\medskip}

% --- Mots-clés ---------------------------------------------------------------
%
% En anglais, en clôture du résumé de la lecture modernisée : le vocabulaire
% moderne sous lequel chercher ce que les feuillets construisent. Le manifeste
% du site (`npm run manifest`) les extrait pour étiqueter la cote entière —
% cette ligne est la source unique des tags, il n'y a pas de fichier à côté.
\newcommand{\keywords}[1]{\par\smallskip\noindent
  {\footnotesize\textsc{Keywords} --- \textit{#1}}\par}

% --- L'appareil critique ----------------------------------------------------

% Le numéro de feuillet, en marge. C'est l'ancre qui permet de régler un
% désaccord sur une formule en regardant *un* feuillet plutôt qu'une chemise
% de six cents. Le site s'en sert aussi pour tourner les pages du fac-similé
% quand on fait défiler la transcription.
\newcommand{\leaf}[1]{%
  \marginnote{\footnotesize\color{gray!70}\textsf{#1}}%
  \label{leaf:#1}\ignorespaces}

% Illisible : on ne devine pas. Un mot inventé qui se lit comme les autres est
% le pire résultat possible.
\newcommand{\ill}{\textcolor{red!60!black}{[\,\ldots\,]}}

% Lecture douteuse : le mot est donné, et signalé comme incertain.
\newcommand{\uncertain}[1]{\uline{#1}}

% Ajout de l'éditeur — une lettre manquante, un mot sous-entendu.
\newcommand{\add}[1]{\textcolor{blue!50!black}{[#1]}}

% Biffé par Grothendieck lui-même. Ce qu'il a rayé fait partie du document :
% une rature dit souvent où la pensée a changé de direction.
\newcommand{\struck}[1]{\sout{#1}}

% Note du transcripteur, sur l'état matériel du feuillet.
\newcommand{\note}[1]{\par\smallskip{\small\color{gray!60!black}\textit{[#1]}}\par\smallskip}

% Note marginale de Grothendieck, distincte du corps du texte.
\newcommand{\marginal}[1]{\par\smallskip\noindent\hspace{1em}%
  {\small\color{gray!60!black}#1}\par\smallskip}

% --- Titre ------------------------------------------------------------------

\AtBeginDocument{%
  \begin{center}
    {\large\bfseries Fonds Alexandre Grothendieck --- cote n\textsuperscript{o}~\grothFolder}\\[2pt]
    {\itshape\grothFoldertitle}\\[4pt]
    {\small feuillets \grothFirst--\grothLast\ (lot \grothBatch)}\\[2pt]
    {\footnotesize\color{gray!60!black}Transcription établie d'après le fac-similé numérisé
     par l'Université de Montpellier.\\
     Les lectures incertaines sont soulignées, les illisibles notées [\,\ldots\,],
     les ajouts entre crochets.}
  \end{center}
  \vspace{1em}}

\endinput


\folder{43}
\batch{3}
\pages{41}{60}
\dating{[à partir de 1961-vers 1968]}
\watermark{Édition de démonstration}
\foldertitle{Dualité des groupes. Jacobiennes généralisées etc : notes et copies de notes manuscrites (s.d.)}

\begin{document}

\page{41}
\note{la page semble continuer les pages 39--40 (descente galoisienne pour le groupe
des classes d'idèles $J_C$ d'une courbe) ; comme au lot précédent, $J$
note son $J$ à double trait, ici et page 42}
\[
\text{\struck{$C^{*}(\Gamma$}}\ \ \varprojlim H^{*}(\pi, J_{C',\alpha}(\bar{k})) \;=\; J_{C'}(k)
\]
\note{sous le premier membre, une accolade renvoie aux deux lignes suivantes}
\[
\begin{array}{ll}
= H^{i}(\pi, \mathbb{Z}) = 0 & \text{si } i > 0 \\
= H^{0}(\pi, J_{C',\alpha}(\bar{k})) = J_{C',\alpha}(k) & \text{\uncertain{si} \ill{} M.L.}
\end{array}
\]
\note{« M.L. » : sans doute Mittag-Leffler ; le mot qui précède est
illisible}
\[
H^{p}\bigl(\varprojlim_{\alpha} H^{q}(\pi, C^{*}(\Gamma, J_{C,\alpha}(\bar{k})))\bigr)
\]
\[
\begin{array}{ll}
\text{\struck{$q > 0$}}\ = H^{q}(\pi, C^{*}(\Gamma, \mathbb{Z})) = 0 & \text{si } q > 0 \\
= H^{0}(\pi, C^{*}(\Gamma, J_{C,\alpha}(\bar{k}))) = C^{*}(\Gamma, J_{C,\alpha}(k)) &
\end{array}
\]
\note{une grande accolade réunit ces lignes en}
\[
C^{*}(\Gamma, J_{C}(k)).
\]

\[
\text{\struck{$\ill{}$}}\ \ltimes E^{1} \to G \qquad \mathrm{Ext}^{1}(J^{0}_{\text{\uncertain{$K$}}}, G)
\]
\note{colonne verticale de flèches descendantes au centre de la page :}
\[
\begin{array}{c}
H^{1}(k, G) \qquad H^{p}(\Gamma, J_{C}(k)) \\
\downarrow \\
H^{1}(K, G) \xleftarrow{\ \sim\ } \mathrm{Ext}^{1}(J_{K}, G) \\
\downarrow \\
\mathrm{Ext}^{1}(J^{0}_{K}, G) \\
\downarrow \\
H^{2}(k, G) = 0
\end{array}
\qquad\qquad \hat{\mathbb{Z}} \to 0
\]
\note{isolés dans la marge gauche : « $G$ \quad $C^{i}(G)$ » et
« $H^{p}(\mathcal{H}$ », inachevé}
\[
\cdots \to H^{0}(k, \Gamma) \to \mathrm{Hom}(J_{K}, \Gamma) \to \mathrm{Hom}(J^{0}_{K}, \Gamma)
\]
\note{sous $H^{0}(k,\Gamma)$ et $\mathrm{Hom}(J_K,\Gamma)$, il
écrit « $\Gamma$ », « $\Gamma$ »}
\[
\boxed{\mathrm{Ext}^{1}(J^{0}, \Gamma) = \mathrm{Hom}}
\]

\page{42}
\[
\mathrm{Ext}^{1}(J, G) \Rightarrow \mathrm{Hom}(J(k), G(k))
\]
\note{le signe entre les deux membres est une flèche double courte ; sous
$\mathrm{Ext}^{1}(J,G)$, une petite flèche descend vers un mot
illisible (\ill{})}
\[
J(k) \times \mathrm{Ext}^{1}(J, G) \longrightarrow
\boxed{H^{1}(k, G) \simeq G(k)}
\]
\note{sous l'encadré, entre parenthèses : « (\uncertain{corps fini}) » ;
sous $J(k)$, un signe d'égalité vertical renvoie à
$H^{0}(k,J)$ ; plus bas, à droite,
$\mathrm{Hom}(J(k), G(k))$, et à gauche
$\mathrm{Hom}(\mathbb{Z}, J)$}
\[
0 \to G \to E \to J \to 0, \qquad \mathbb{Z} \to J, \qquad J'
\]
\note{la flèche $\mathbb{Z} \to J$ monte verticalement sous
$J$ ; un double trait en marche d'escalier sépare cette suite de
« $J'$ », écrit en haut à droite}
\[
\pi_{1}(J) - \pi_{1}(\mathbb{Z})
\]

\begin{tikzcd}[column sep=small, row sep=small, nodes={font=\scriptsize}]
0 \arrow[r] & \mathrm{Ext}^{1}(\mathbb{Z}, G) \arrow[r] \arrow[d, "\wr"] & \mathrm{Ext}^{1}(J, G) \arrow[r] \arrow[d] & \mathrm{Ext}^{1}(J^{0}, G) \arrow[r] \arrow[d, "\wr"] & 0 \\
0 \arrow[r] & \mathrm{Hom}(\mathbb{Z}, H^{1}(G)) \arrow[r] & \mathrm{Hom}(J(k), H^{1}(G)) \arrow[r] & \mathrm{Hom}(J^{0}(k), H^{1}(G)) \arrow[r] & 0
\end{tikzcd}

\note{les deux flèches verticales extrêmes portent un signe que nous
lisons comme $\wr$ (isomorphisme) ; il pourrait aussi s'agir d'un « 2 »}

\page{44}
\marginal{Adèles pour plusieurs dimensions}
\note{encadré de sa main, en haut à droite de la page ; la page est une
suite de formules sans phrase}
\[
\prod_{x \in X} \hat{\mathcal{O}}_{x}
\]
\[
\prod^{\mathrm{local}}_{x_1 \in X_1} \hat{\mathcal{O}}_{x_1} \otimes_{\underline{\mathcal{O}}_{x_1}} K
\;=\; \prod_{\mathrm{local}} K_{D}
\]
\[
\prod_{x_1} \text{\struck{$\ill{}$}} \prod^{\mathrm{local}}_{x_0 \in \overline{x_1}} \mathcal{O}_{[x_0, x_1]}
\]
\[
\prod_{x_0 x_1 (x_2)} \mathcal{O}_{[x_0 x_1 x_2]}
\;\longrightarrow\;
\mathcal{O}_{[x_0 x_1]} \otimes_{\mathcal{O}_{x_1}} K,
\qquad
\mathcal{O}_{[x_0 x_\nu]} = \hat{\underline{\mathcal{O}}}_{x_0} \otimes_{\underline{\mathcal{O}}_{x_0}} K
\]
\note{sur la page la flèche est verticale ; à côté, un croquis : deux
courbes qui se coupent en un point $x_0$, l'une allant vers $x_1$}
\[
K \longrightarrow I_{x_1} \longrightarrow I_{x_0}
\]
\note{sous $K$, un signe d'égalité vertical ; sous la première flèche,
une accolade renvoie à $\hat{\mathcal{O}}_{x_1} \otimes_{\underline{\mathcal{O}}_{x_1}} K$,
sous la seconde à $\mathcal{O}_{[x_0, x_1]}$ ; en dessous, alignés :
$\hat{\underline{\mathcal{O}}}_{x_{\uncertain{2}}}$,
$\hat{\underline{\mathcal{O}}}_{x_{\uncertain{1}}}$,
$\hat{\underline{\mathcal{O}}}_{x_{0}}$ ; les indices sont mal formés}
\[
\hat{\underline{\mathcal{O}}}_{x_0} \quad \hat{\mathcal{O}}_{x_1} \quad \hat{\mathcal{O}}_{x_2}
\qquad\qquad
\mathfrak{m}_{0} \quad \mathfrak{m}_{1} \quad \mathfrak{m}_{2}
\]
\[
\prod_{x \in X} \mathcal{O}_{[x]}, \qquad
\prod^{\mathrm{local}}_{x_1 \in X_1} \mathcal{O}_{[x_1, x_2]}, \qquad \mathcal{O}.
\]
\note{l'indice $[x_1, x_2]$ est bien écrit ainsi}

\begin{tikzcd}[column sep=small, row sep=small, nodes={font=\scriptsize}]
\prod_{x_0} \mathcal{O}_{[x_0]} \arrow[d, "\mathrm{localisant}"] & & \\
\prod_{x_1} \prod^{\mathrm{local}}_{x_0} \mathcal{O}_{[x_0, x_1]} \arrow[d, "\mathrm{localisant}"] & \prod_{x_1} \mathcal{O}_{[x_1]} \arrow[l] \arrow[d, "\mathrm{localisant}"] & \\
\prod^{\mathrm{vert.}}_{x_0} \prod^{\mathrm{vert.}}_{x_1} \mathcal{O}_{[x_0, x_1, x_2]} & \prod^{\mathrm{local}}_{(x_1, x_2)} \mathcal{O}_{[x_1 x_2]} \arrow[l] & \prod_{x_2} \mathcal{O}_{[x_2]} = K \arrow[l]
\end{tikzcd}

\note{diagramme redessiné : sur la page, les trois étages sont décalés en
escalier et les indices $x_0$, $x_1$ des produits sont soulignés ; à gauche
de ce diagramme, séparé par un trait oblique, il répète
$\prod_{x_1} \prod^{\mathrm{local}}_{x_0} \mathcal{O}_{[x_0, x_1]}$ ; « vert. »
se lit peut-être « vertical »}
\[
\underline{\mathcal{O}}_{[x_1, x_2]}, \qquad \hat{\mathcal{O}}_{x_0}
\]
\note{à côté de $\hat{\mathcal{O}}_{x_0}$, un croquis : une courbe passant
par $x_0$ et entourée d'un ovale ; plus bas, une ligne biffée
« $\struck{\ill{} \in K_i, c}$ »}
\[
\mathcal{O}_{[x_0 x_1]} \longrightarrow \prod_{x_0 x_1} \mathcal{O}_{[x_0 x_1 x_2]}
\]
\[
\text{\struck{$\mathcal{A}_{i_0 \cdots i_{p-1}}$}}\quad
\mathcal{A}_{i_0 \cdots i_p} \qquad \mathcal{A}_{i_1 \cdots i_{p-1}}
\]
\[
x_{i_0}\ x_{i_1} \cdots x_{i_p}, \qquad i_0 < i_1 < \cdots < i_p
\]

\page{45}
\[
\mathcal{A}^{*} = \prod_{x \in X} \mathcal{A}^{*}\pi_{x}
\]
\[
\mathcal{A}^{*}\pi_{x} = \varprojlim_{n} (\mathcal{A}^{*}\pi_{x}) /
\]
\note{la formule s'interrompt après la barre de quotient ; le reste de
la page est blanc}

\page{46}
\[
\varprojlim_{i}\ \mathcal{A}^{0} \otimes_{A} \mathcal{A}^{0}/\mathfrak{m}_{i}
\qquad\qquad
\prod \hat{\underline{\mathcal{O}}}_{x} \otimes_{A} K
\]
\[
\mathcal{A}^{0} \otimes_{A} \mathcal{A}^{0}/\mathfrak{m}_{i}
\;\subset\; \prod_{x \in X} \hat{\underline{\mathcal{O}}}_{x} \otimes_{A} \mathcal{O}_{a}/\mathfrak{m}_{a}^{n}
\;=\; \prod_{x \leq a} \hat{\mathcal{O}}_{x} \otimes \mathcal{O}_{a}/\mathfrak{m}_{a}^{n}
\]
\note{sous le second produit, un indice biffé et illisible, remplacé par
$x \leq a$}
\[
\prod_{a} \varprojlim_{n} \Bigl(\prod_{x \leq a} \hat{\mathcal{O}}_{x} \otimes \mathcal{O}_{a}/\mathfrak{m}_{a}^{n}\Bigr)
\]
\[
= \text{\struck{$a \in X_{\ill{}}$}}\ \prod_{a} \varprojlim_{n} \Bigl[\bigl(\hat{\mathcal{O}}_{a} \otimes \underline{\mathcal{O}}_{a}/\mathfrak{m}_{a}^{n}\bigr) \times \Bigl(\bigl(\prod_{x < a} \hat{\mathcal{O}}_{x}\bigr) \otimes \mathcal{O}_{a}/\mathfrak{m}_{a}^{n}\Bigr)\Bigr]
\]
\[
= \prod_{a} \hat{\mathcal{O}}_{a} \times \prod_{i} \prod_{a \in X_i} \varprojlim_{n} \Bigl(\bigl(\prod_{x < a} \hat{\mathcal{O}}_{x}\bigr) \otimes \underline{\mathcal{O}}_{a}/\mathfrak{m}_{a}^{n}\Bigr)
\]
\note{une accolade sous le dernier facteur renvoie à trois cas, écrits en
colonne ; à gauche, deux lignes biffées et illisibles}
\[
\begin{array}{ll}
i = 0 & \text{\ill{}} \\
i = 1 & \bigl(\prod_{x < a} \hat{\mathcal{O}}_{x}\bigr) \otimes \mathcal{O}_{a}/\mathfrak{m}_{a}^{n} \\
i = 2 & \bigl(\prod_{x < a} \hat{\mathcal{O}}_{x}\bigr) \otimes K
\end{array}
\]
\[
\mathcal{A}^{n} = \text{\struck{$\mathcal{A}$}}\ \varprojlim_{\alpha} \mathcal{A}^{n-1} \otimes_{A} \mathcal{A}^{0}/\mathfrak{m}_{\alpha}
= \prod_{a \in X} \varprojlim_{k} \mathcal{A}^{n-1} \otimes_{A} \mathcal{O}_{a}/\mathfrak{m}_{a}^{k}
\]
\[
\cap \qquad \prod_{x_0 \leq x_1 \leq \cdots \leq x_n} \mathcal{A}_{[x_0, \ldots, x_n]}
\]
\note{le signe $\cap$ est écrit sous la limite, comme une inclusion
verticale ; à droite, souligné : « l'inclusion étant fidèlement plate »}
\[
\mathcal{A}^{n+1} = \prod_{a \in X} \varprojlim_{k} \mathcal{A}^{n} \otimes_{A} \mathcal{O}_{a}/\mathfrak{m}_{a}^{k}
\;\subset\; \prod_{a \in X} \varprojlim_{k} \Bigl(\prod_{x_0 \leq \cdots \leq x_n} \mathcal{A}_{[x_0 \ldots x_n]}\Bigr) \otimes \mathcal{O}_{a}/\mathfrak{m}_{a}^{k}
\]
\[
\subset \prod_{x_0 \leq \cdots \leq x_n} \mathcal{A}_{[x_0 \ldots x_n]} \otimes \underline{\mathcal{O}}_{\text{\uncertain{$a$}}}\ \ill{}
\qquad \mathcal{A}_{[x_0, \ldots, x_n, a]}
\]
\note{une flèche double horizontale relie le produit entre parenthèses à
la dernière ligne ; celle-ci sort de la page à droite}

\page{47}
\note{page de calculs épars, sans phrase suivie hormis une glose ; nous
les donnons colonne par colonne, de gauche à droite ; en haut à gauche, un
croquis : une courbe et, en un point marqué, deux petits cercles}
\[
\begin{array}{c}
\mathcal{A}_{i_0 \cdots i_{p-1}} \\
\downarrow \\
\mathcal{A}_{i_0 \cdots i_p} \longleftarrow \mathcal{A}_{i_p}
\end{array}
\]
\[
\mathcal{A}_{\ell_1, \ell_2, \ldots, \ell_p} \qquad \mathcal{A}_{\ell_{p+1}} \qquad \mathcal{A}_{\ell}
\qquad \mathcal{A}_{1} \qquad \text{\struck{$\mathcal{O}_{[}$}}\ \mathcal{A}_{0} \otimes
\]
\note{la colonne centrale :}
\[
\prod \hat{\underline{\mathcal{O}}}_{x} \otimes_{\mathcal{O}_x} K,
\qquad
\hat{\mathcal{O}}_{x} \otimes_{\underline{\mathcal{O}}_x} \text{\struck{$\ill{}$}}\ L_{n}
\]
\[
\hat{\mathcal{O}}_{x} \otimes_{\underline{\mathcal{O}}_x} (\mathfrak{p}^{n}/\mathfrak{p}^{n+1})
\to \hat{\mathcal{O}}_{x} \otimes L_{n+1} \to \hat{\underline{\mathcal{O}}}_{x} \otimes L_{n}
\]
\[
\hat{\mathcal{O}}_{x}/\mathfrak{p}^{n+1} \longrightarrow \hat{\mathcal{O}}_{x}/\mathfrak{p}^{n}
\]
\note{des flèches montantes relient ces deux lignes ; elles sont marquées
d'astérisques, ainsi que les termes de la seconde ligne (groupes des
unités ?)}
\[
\hat{\mathcal{O}}_{x} \otimes_{\mathcal{O}_{\mathfrak{p}}} \mathfrak{p}^{n}\mathcal{O}_{\mathfrak{p}}/\mathfrak{p}^{n+1}\mathcal{O}_{\mathfrak{p}}
\to \hat{\mathcal{O}}_{x} \otimes_{\mathcal{O}_{\mathfrak{p}}} \mathcal{O}_{\mathfrak{p}}/\mathfrak{p}^{n}\mathcal{O}_{\mathfrak{p}}
\to \hat{\mathcal{O}}_{x} \otimes_{\mathcal{O}_{\mathfrak{p}}} L
\]
\note{dans le premier terme, un symbole est biffé après le produit
tensoriel}
\[
\hat{\mathcal{O}}_{x} \to \hat{\mathcal{O}}_{x}\, S^{-1} \qquad \mathfrak{p} \quad \mathfrak{p}^{2}
\]
\[
\mathcal{A}_{[\mathfrak{p}_0, \ldots, \mathfrak{p}_n]} \otimes_{A} A_{\mathfrak{p}_{n+1}}/\mathfrak{p}_{n+1}^{\text{\uncertain{$n$}}} A_{\mathfrak{p}_{n+1}}
\qquad
A_{[\mathfrak{p}_0]} \mathbin{\hat{\otimes}_{A}} A_{[\mathfrak{p}_1]}
\qquad
\hat{\underline{\mathcal{O}}}_{x} \otimes K
\]
\[
A_{[x]} \otimes \qquad \varprojlim \hat{A} \otimes_{A} A/\mathfrak{m}_{i} \qquad \varprojlim A
\]
\note{la colonne de droite : un calcul local en deux variables, avec un
croquis (une courbe coupée par une autre en un point, « $s = 0$ »)}
\[
\begin{array}{c}
k[[s,t]] \\
\uparrow \\
k[s,t]_{\mathfrak{m}} \longrightarrow k[s,t]_{(s)} \\
\uparrow \\
k[t]
\end{array}
\]
\note{les flèches sont annotées verticalement « $x_0 x_1$ », « $x_1 x_2$ »,
« $x_0 x_2$ » ; à côté, écrit en travers et en partie biffé :
$(\hat{\mathcal{O}}_{x}^{*} \otimes K^{*})$}
\[
k[[t]][s]\,(s^{-1}) \qquad k[[s]][t^{-1}]
\]
Localisé en les polynômes en $s, t$ qui ne sont \uncertain{pas} multiples de $s$
\[
a_{0}(t) + a_{1}(t)\,s + \cdots + a_{n-1}(t)\,s^{n-1}
\]
\[
a_{0}(t)\,[1 + a_{0}(t)^{-1} a_{1}(t)\,s + \cdots
\]
\[
\text{\struck{$a_0$}}\ k((t))[s]/(s^{n}) \longleftarrow k[[s,t]]
\]
\note{la flèche est verticale et monte vers $k((t))[s]/(s^{n})$}
\[
\underline{\underline{a_{0}(t)}} + \underline{\underline{a_{1}(t)}}\,s + \cdots + \underline{\underline{a_{n-1}(t)}}\,s^{n-1}
\]
\[
k((t))[[s]] \qquad k[[t]][[s]] \qquad k((t))/k[[t]]\ \ [[s]]
\]

\page{48}
\note{feuillet écrit dans le sens de la largeur ; calculs épars, sans
phrase ; nous les donnons par blocs, de gauche à droite et de haut en bas}
\[
\prod_{x \in X_0} G(\hat{\mathcal{O}}_{x}), \qquad
\prod_{D} \prod^{\mathrm{local}}_{x \in |D|} G(\hat{\mathcal{O}}_{D,x}), \qquad
\prod_{\text{\struck{$\ill{}$}}\ D \in X_1} G(\hat{\mathcal{O}}_{D}), \qquad
\prod_{\mathrm{local}} G(\hat{K}_{D})
\]

\begin{tikzcd}[column sep=small, row sep=small, nodes={font=\scriptsize}]
\prod_{x} \hat{\underline{\mathcal{O}}}_{x}^{*} \arrow[d] & & \\
\prod_{D} \prod^{\mathrm{local}}_{x \in |D|} \hat{\mathcal{O}}_{D,x}^{*} \arrow[d] & \prod_{D} \hat{\mathcal{O}}_{D}^{*} \arrow[l] \arrow[d] & 0 \arrow[l, dashed] \arrow[d, dashed] \\
\prod_{x} \prod_{|D| \ni x} (\hat{\mathcal{O}}_{D,x} \otimes_{\underline{\mathcal{O}}_{D}} K)^{*} & \prod_{\mathrm{local}} \hat{K}_{D}^{*} \arrow[l] & K^{*} \arrow[l]
\end{tikzcd}

\note{diagramme redessiné ; la dernière colonne est tracée en pointillé
depuis le « $0$ » ; à gauche de la dernière ligne, séparé par un trait
oblique, « $G(K)$ » ; sous le second produit, « $D$ » est surchargé}
\[
\mathrm{Hom}\bigl(K, \textstyle\coprod_{x} \text{\struck{$\ill{}$}}\, I_{x}\bigr)
\]
\[
K \to I_{D} \to I_{x}
\]
\note{sous la première flèche, $\hat{K}_{D}^{*}$ ; sous la seconde,
barré d'un trait, $\hat{\mathcal{O}}_{D,x}^{*}$}
\[
\prod_{D} \hat{\mathcal{O}}_{D}^{*} \longleftarrow \Sigma\, \underline{\mathcal{O}}_{D}, \qquad K^{*}, \qquad \mathcal{O}_{D}
\]

\begin{tikzcd}[column sep=small, row sep=small]
\hat{\underline{\mathcal{O}}}_{D_x} & \hat{\mathcal{O}}_{D} \arrow[l] \arrow[d] & \underline{\mathcal{O}}_{D} \arrow[l] \arrow[d] \\
 & \hat{K}_{D} & K^{*} \arrow[l]
\end{tikzcd}

\note{sous ce carré, à gauche, une rature pleine}

\note{colonne de droite, avec un croquis en haut (une courbe, et une
autre traversant un disque hachuré) :}
\[
\begin{array}{c}
\hat{\underline{\mathcal{O}}}_{x} \leftarrow \mathcal{O}_{x} \\
\downarrow \\
\hat{\underline{\mathcal{O}}}_{x} \otimes_{\underline{\mathcal{O}}_x} \underline{\mathcal{O}}_{D} \leftarrow \mathcal{O}_{D} \\
\downarrow \\
\underline{\mathcal{O}}_{\text{\uncertain{$\mathfrak{q}$}}}
\end{array}
\qquad\qquad
\begin{array}{ccc}
\hat{\mathcal{O}}_{D} & \leftarrow & \mathcal{O}_{D} \\
\downarrow & & \downarrow \\
\hat{\underline{\mathcal{O}}}_{D,0} & \leftarrow & \mathcal{O}_{0}
\end{array}
\]
\note{le premier membre de la seconde ligne de la colonne gauche porte un
chapeau commun et un trait biffé}

\note{colonne de gauche, en bas :}
\[
\begin{array}{l}
A_{[\mathfrak{p}_1, \ldots, \mathfrak{p}_n]} \\
\uparrow \ \text{hom.\ local} \\
A_{[\mathfrak{p}_n]}\ \ \text{\struck{$\leftarrow A_{[\mathfrak{p}_{n-1}]}/\mathfrak{p}_{n-1}^{k} A_{[\mathfrak{p}_{n-1}]}$}} \\
\uparrow \ \text{hom.\ local} \\
A_{\mathfrak{p}_n} \longrightarrow A_{\mathfrak{p}_{n+1}}/\mathfrak{p}_{n+1}^{k} A_{\mathfrak{p}_{n+1}}
\end{array}
\]
\[
\text{\struck{$A_{[\mathfrak{p}_1]} \leftarrow A_{[\mathfrak{p}_2]} \to A_{[\mathfrak{p}_1]}$}}
\qquad
A_{[\mathfrak{p}_1]}, \qquad
A_{\mathfrak{p}_1} \longrightarrow \text{\struck{$A$}}\,A_{\mathfrak{p}_2}/\mathfrak{p}_2^{k} A_{\mathfrak{p}_2}, \qquad
\underline{A}
\]
\note{les deux premières lignes de ce dernier groupe, où la flèche
$A_{[\mathfrak{p}_2]} \to A_{[\mathfrak{p}_1]}$ monte, sont biffées de
traits obliques}

\page{49}
\emph{Détermination des} $\lbrace D_{\mathfrak{p}}(M) \longrightarrow K_{\mathfrak{q}} \rbrace$ \quad $(\mathfrak{q} \supset \mathfrak{p})$
\note{le mot après « des » est une abréviation soulignée que nous ne
lisons pas avec sûreté (« hom. » ?) ; l'accolade ouvrante qui la suit peut
aussi être une parenthèse}

On peut supposer $A$ local d'idéal maximal $\mathfrak{m} = \mathfrak{q}$,
$\mathfrak{p}$ sera un idéal premier quelconque. \struck{On a}
$\mathrm{Hom}(D_{\mathfrak{p}}(M), K_{\mathfrak{q}})$ est un foncteur
\emph{\uncertain{exact}} en $M$,\marginal{\uncertain{covariant}} donc de la
forme $M \otimes K_{\mathfrak{p},\mathfrak{q}}$, où
$\mathbf{K}_{\mathfrak{p},\mathfrak{q}}$ est un $A$-module
\struck{\ill{}} \emph{plat}, savoir
$\mathrm{Hom}(K_{\mathfrak{p}}, K_{\mathfrak{q}})$. Si
$\mathfrak{q} = \mathfrak{p}$, on obtient
$\hat{A}_{\mathfrak{p}} = \widehat{A_{\mathfrak{q}}}$.
\note{la seconde occurrence de $K_{\mathfrak{p},\mathfrak{q}}$ est
écrite en gras (repassée) ; la dernière égalité est lue telle quelle,
malgré $\mathfrak{q} = \mathfrak{p}$}

\page{50}
$A$ anneau \uncertain{régulier}, de \uncertain{dim.}\ finie, \struck{\ill{}}
\[
\text{\struck{$D(A) = \coprod_{\mathfrak{p}} D_{\mathfrak{p}}(A)$}}
\qquad
\text{\struck{$DD(A) = \prod_{\mathfrak{p}} \mathrm{Hom}(I_{\mathfrak{p}},$}}
\]
\note{ces deux lignes sont barrées d'une croix ; sous $D_{\mathfrak{p}}(A)$,
une accolade renvoie à $I_{\mathfrak{p}}$}
\[
D(M) = \mathrm{Hom}(M, K_{*}) = \coprod_{\mathfrak{p}} \mathrm{Hom}(M, K_{\mathfrak{p}})
\]
\note{sous le dernier terme, une accolade : « $\mathfrak{p}$-primaire »,
« $D_{\mathfrak{p}}(M)$ »}
\[
DD(M) = \prod_{\mathfrak{p}} \mathrm{Hom}(D_{\mathfrak{p}}(M), K)
\]
\note{$D_{\mathfrak{p}}(M)$ est surchargé sur une première lettre}
\[
D_{\mathfrak{p}}(M) = \varinjlim_{i} U_{i}
\]
les $U_i$ \uncertain{sont} de longueur finie sur $A_{\mathfrak{p}}$
\[
\mathrm{Hom}(D_{\mathfrak{p}}(M), K) = \varprojlim_{i} \mathrm{Hom}(U_{i}, K)
\]
\[
\mathrm{Hom}(U_{i}, K_{*}) = \coprod_{\mathfrak{q}} \mathrm{Hom}(U_{i}, K_{\mathfrak{q}})
\]
\[
\mathrm{Hom}(U_{i}, K_{\mathfrak{q}}) = \varinjlim_{j} \mathrm{Hom}(U_{i}, K_{\mathfrak{q},j})
\]
\note{sous $K_{\mathfrak{q},j}$, une accolade : « $V_j$ de longueur finie
sur $A_{\mathfrak{q}}$ » ; dans la marge gauche, un croquis : deux courbes
qui se coupent en un point marqué}

Or $\mathrm{Hom}(U_i, V_j) = 0$ si $\mathfrak{q} \not\supset \mathfrak{p}$ ;
donc \ill{} \ill{} $\mathrm{Hom}(U_i, K_{\mathfrak{q}}) = 0$ si
$\mathfrak{q} \not\supset \mathfrak{p}$, d'où
\[
\mathrm{Hom}(U_i, K_{*}) = \coprod_{\mathfrak{q} \supset \mathfrak{p}} \mathrm{Hom}(U_i, K_{\mathfrak{q}})
\]
\note{sous le coproduit, l'indice $\mathfrak{q} \supset \mathfrak{p}$ est
écrit au-dessus d'un indice raturé ; de même dans les deux formules
suivantes}
\[
DD(M) = \prod_{\mathfrak{p}} \mathrm{Hom}\Bigl(D_{\mathfrak{p}}(M), \coprod_{\mathfrak{q} \supset \mathfrak{p}} K_{\mathfrak{q}}\Bigr)
= \prod_{\mathfrak{p}} \prod^{\mathrm{local}}_{\mathfrak{q} \supset \mathfrak{p}} \mathrm{Hom}(D_{\mathfrak{p}}(M), K_{\mathfrak{q}})
\]

\page{51}

\begin{tikzcd}[column sep=small, row sep=small]
0 & \hat{\mathcal{O}}_{\mathfrak{q}_1} \arrow[l] \arrow[d] & \mathcal{O}_{\mathfrak{q}_1} \arrow[l] \arrow[d] \\
 & \hat{\mathcal{O}}_{\mathfrak{q}_1 \mathfrak{q}_0} & \mathcal{O}_{\mathfrak{q}_0} \arrow[l]
\end{tikzcd}

\note{au milieu du carré, le signe $\otimes$ ; en haut, au centre, un
croquis (deux arcs issus d'un point $x$) et, à droite, un segment à trois
points marqués ; plus bas : « $x$ \quad $\supset$ \quad \uncertain{finis} »
et, entouré d'un ovale, $\hat{\underline{\mathcal{O}}}^{*}_{x} \otimes_{\underline{\mathcal{O}}_x} K$}
\[
A_{[\mathfrak{p} \supset \mathfrak{q}]}
\]
\[
\begin{array}{c}
A_{[\mathfrak{p}]} = \widehat{A_{\mathfrak{p}}} \\
\uparrow \\
A_{\mathfrak{p}} \longrightarrow A_{\mathfrak{q}}/\mathfrak{q}^{n} A_{\mathfrak{q}}
\end{array}
\]
\[
A_{[\mathfrak{p}]} = \widehat{A_{\mathfrak{p}}} \longleftarrow A_{\mathfrak{p}}
\]
\[
\begin{aligned}
A_{[\mathfrak{p}, \mathfrak{q}]} &= \varprojlim_{n} A_{[\mathfrak{p}]} \otimes_{A} A_{\mathfrak{q}}/\mathfrak{q}^{n} A_{\mathfrak{q}} \\
&= A_{[\mathfrak{p}]} \otimes_{A_{\mathfrak{p}}} A_{\mathfrak{q}}/\mathfrak{q}^{n} A_{\mathfrak{q}} \\
&= A_{[\mathfrak{p}]}/\mathfrak{q}^{n} A_{[\mathfrak{p}]} \otimes_{\text{\struck{$A_{\mathfrak{q}}$}}\, A_{\mathfrak{p}}} A_{\mathfrak{q}}/\mathfrak{q}^{n} A_{\mathfrak{q}} \\
&= \text{\struck{$(A/\mathfrak{q}^{n} A$}} \\
&= A_{[\mathfrak{p}]}/\mathfrak{q}^{n} A_{[\mathfrak{p}]} \otimes_{A_{\mathfrak{p}}} A_{\mathfrak{q}} \\
&= A_{[\mathfrak{p}]}/\mathfrak{q}^{n} A_{[\mathfrak{p}]} \otimes_{A} A_{\mathfrak{q}}
\end{aligned}
\]
\note{les égalités sont écrites verticalement ; une accolade sous la
limite projective mène à la deuxième ligne, qui ne porte plus de limite ;
la dernière ligne est soulignée et une flèche en part vers la formule
suivante}
\[
\widehat{A_{\mathfrak{p}}/\mathfrak{q}^{n} A_{\mathfrak{p}}} = (A/\mathfrak{q}^{n} A)_{[\mathfrak{p}]}
\qquad\qquad
\text{\struck{$\ill{}\ (A/\mathfrak{q}^{n} A)_{[\mathfrak{p}]}$}}
\]
\marginal{\uncertain{on peut y remplacer} $\mathfrak{q}^{n}$ par $\mathfrak{q}^{(n)}$}
\note{note dans une bulle à gauche, reliée par un trait sinueux à la
formule précédente}
\marginal{\ill{} \ill{} \ill{}, \ill{} \ill{} \ill{} \ill{} \ill{} \ill{}
\ill{} \ill{} \ill{} \ill{} \ill{} \ill{} \ill{} de
$\widehat{A_{\mathfrak{p}}/\mathfrak{q}A_{\mathfrak{p}}} \otimes_{A_{\mathfrak{p}}/\mathfrak{q}A_{\mathfrak{p}}}$ \uncertain{$1$}}
\note{six lignes serrées dans la marge droite, séparées du calcul par une
grande parenthèse ; seule la formule finale se lit}

\page{52}
\[
\mathfrak{p} \longrightarrow \mathfrak{m}, \qquad \mathfrak{p}' \quad \mathfrak{m}'
\]
\note{en haut de page, deux schémas raturés : $\mathfrak{p}$ et
$\mathfrak{m}$ reliés par une flèche, surmontés de $\mathfrak{p}'$ et
$\mathfrak{m}'$ surchargés, le tout souligné d'un crochet ; à droite, un
croquis en escalier (des cases empilées dans l'angle d'un cadre, une flèche
descendante et une flèche vers la gauche)}
\[
\mathrm{Hom}(U_{\mathfrak{p}}, K_{\mathfrak{m}}) = A_{\mathfrak{m},\mathfrak{p}}
\]
\[
\mathfrak{p} \subset \mathfrak{p}' \qquad \mathfrak{q}' \subset \mathfrak{q} = \mathfrak{m}
\qquad\qquad
A_{\mathfrak{m}}/\mathfrak{p}^{n} A_{\mathfrak{m}} = \mathbf{A}
\]

\begin{tikzcd}[column sep=small, row sep=small]
A_{\mathfrak{q},\mathfrak{p}''} \arrow[d] & & \\
A_{\mathfrak{q},\mathfrak{p}'} \arrow[d] & A_{\mathfrak{q}',\mathfrak{p}'} \arrow[l] \arrow[d] & \\
A_{\mathfrak{q},\mathfrak{p}} & A_{\mathfrak{q}',\mathfrak{p}} \arrow[l] & A_{\mathfrak{q}'',\mathfrak{p}} \arrow[l]
\end{tikzcd}

\note{dans les termes de gauche, l'indice est surchargé (« $\mathfrak{q}$ »
écrit sur une autre lettre)}
\[
\mathfrak{p} \subsetneq \mathfrak{p}' \subsetneq \mathfrak{p}'' \subset \mathfrak{q}' \subsetneq \mathfrak{q},
\qquad
\mathfrak{p} \subsetneq \mathfrak{p}' \subseteq \mathfrak{q}' \subsetneq \mathfrak{q},
\qquad
\mathfrak{p} \subset \mathfrak{q}'' \subsetneq \mathfrak{q}' \subsetneq \mathfrak{q}
\]
\note{ce ne sont pas des formules linéaires : les idéaux sont disposés en
losange, reliés par des traits d'inclusion ; nous en donnons les trois
chaînes}
\[
\mathrm{Hom}(A_{\mathfrak{p}}, K_{\mathfrak{q}}), \qquad
\mathrm{Hom}(K_{*}, K_{*})
\]
\note{sous $\mathrm{Hom}(K_*,K_*)$, une flèche à deux pointes ; entre ses
arguments, un petit indice illisible}
\[
A_{\mathfrak{q},\mathfrak{p}} = \Bigl(\coprod_{\mathfrak{p}'} A_{\mathfrak{q},\mathfrak{p}'}\Bigr) \amalg \Bigl(\coprod_{\mathfrak{q}'} A_{\mathfrak{q}',\mathfrak{p}}\Bigr) \Big/
\]
\note{la formule s'arrête sur la barre de quotient ; le second coproduit
est surchargé}

$I$ ensemble ordonné, tel que
(i) \uncertain{deux} \uncertain{objets} \struck{\ill{}} $x, y$ tels que
$x \leq y$ \uncertain{peuvent} être \uncertain{joints} par une
\uncertain{chaîne} \uncertain{saturée} ;
(ii) Deux \uncertain{chaînes} \uncertain{saturées} ont \uncertain{même}
\uncertain{longueur}.
\struck{Si $x, y$} \uncertain{On} \uncertain{suppose} \uncertain{défini}
$A_{x,y} \in \mathcal{C}$ si $x \leq y$ et $d(x,y) \leq 1$, \ill{} \ill{}
\ill{} \struck{\ill{}, \ill{}} \ill{} : \ill{} \ill{} \ill{},
\note{avant « $x \leq y$ », un « $d(x,y)$ » biffé ; les deux dernières
lignes ne se lisent pas au-delà de quelques lettres}

\page{53}
\[
\begin{aligned}
A_{[\mathfrak{p},\mathfrak{q}]} &= \varprojlim A_{[\mathfrak{p}]} \otimes_{A_{\mathfrak{p}}} A_{\mathfrak{q}}/\mathfrak{q}^{n} A_{\mathfrak{q}} \\
&= \varprojlim A_{[\mathfrak{p}]}/\mathfrak{q}^{n} A_{[\mathfrak{p}]} \otimes_{A_{\mathfrak{p}}} A_{\mathfrak{q}}/\mathfrak{q}^{n} A_{\mathfrak{q}}
\end{aligned}
\]
\note{dans la seconde ligne, l'indice du second $A_{[\mathfrak{p}]}$ est
surchargé}

$A_{[\mathfrak{p},\mathfrak{q},\mathfrak{r}]}$ \quad etc.

$X$ \emph{ensemble ordonné}

$S_X$ ensemble simplicial associé, ayant les $n$-simplexes \uncertain{sur}
les chaînes
\[
x_0 \leq x_1 \leq \cdots \leq x_n
\]
\note{sous la chaîne, une flèche vers la gauche}
[i.e.\ les applications \uncertain{croissantes} de $[0,n] = \Delta_n$ dans
$X$]

À tout $\sigma \in S_X$ est associé un $A_\sigma$.

À toute opération \struck{face} simpliciale
\note{« simpliciale » est écrit au-dessus de « face » biffé}
$u : \Delta_m \to \Delta_n$ (\emph{croissante}) est associé un hom.
\[
\text{\struck{Ah}}\quad S_X^{n} \xrightarrow{\ u^{*}\ } S_X^{m}
\]
\note{le mot biffé devant $S_X^n$ n'est pas sûr}
et si $\sigma \in S_X^n$, d'où $u^{*}(\sigma) \in S_X^m$, on a un
morphisme
\[
\varphi(\sigma, u) : A_\sigma \longleftarrow A_{u^{*}(\sigma)},
\]
\uncertain{De} \uncertain{façon} \uncertain{transitive}.

\page{54}
Ainsi, les $A_\sigma$ forment un \emph{faisceau} sur l'ensemble
simplicial $S_X$. \struck{On rappelle} \struck{les faisceaux \ill{}}\note{le
mot biffé après « faisceaux », souligné, se lit peut-être
« cohérents »} [\uncertain{les} \ill{} \ill{} \ill{} \ill{} un
objet simplicial dans la catégorie des ensembles
\uncertain{ordonnés}\note{« ordonnés » est écrit au-dessus de deux mots
biffés et illisibles}].
\struck{\ill{} \ill{} \ill{}} Ici, \struck{\ill{}} \ill{} \ill{} de
dégénérescences \uncertain{compliquées}
\struck{\ill{} \ill{} \ill{} \ill{}}
\ill{} \ill{} \ill{} \ill{} \ill{} \ill{} \ill{} \ill{} \ill{}
\uncertain{actuel}, il y a lieu de \emph{négliger} \ill{} \ill{} qui
\uncertain{proviennent} des simplexes de la forme
\struck{$\partial_i \sigma$, \ill{} \ill{}}
\struck{\ill{} \ill{} \uncertain{simplexes} \ill{}} \emph{dégénérés}
\struck{[i.e.\ des chaînes}
$\partial_i \sigma$, \ill{} $\partial_i$ \ill{} \ill{} opération
\uncertain{face} \ill{} \ill{} \ill{} \ill{} \ill{} une opération
« de \uncertain{coin} », $\sigma$ \ill{} \emph{dégénéré}
[\uncertain{cela} \uncertain{revient} à \uncertain{ne} \uncertain{retenir}
\uncertain{que} les chaînes \uncertain{strictes} \struck{\ill{}}, et
\struck{\ill{}} \ill{} \ill{} \ill{} \ill{} \uncertain{introduire}
\ill{} \ill{} \ill{} \uncertain{équivalences}).
\note{prose très rapide, surchargée de ratures et d'interlignes ; une
flèche venue de la marge gauche renvoie à l'interligne du cinquième
paragraphe ; nous ne donnons que ce qui se lit}

\note{à gauche, des croquis de segments subdivisés, avec les légendes :}
\[
(x_0 - \cdots - x_n), \qquad
(x_0 - \overset{\downarrow}{x_p} - x_n), \qquad
x_0\ x_1\ x_2, \qquad x_0 - x_2
\]

\page{55}
\marginal{[sans \uncertain{vérif.}\ des opérations de dégénérescence)}
\note{en haut à droite, entre un crochet ouvrant et une parenthèse
fermante}

$K^{*}$ \struck{\ill{}} objet cosimplicial dans une catégorie abélienne
$\mathcal{C}$.
\note{« $\mathcal{C}$ » est ajouté au-dessus de la fin de la ligne}
C'est donc un \struck{objet} $\mathcal{C}$-module \uncertain{gradué}
\uncertain{muni} \ill{} \ill{} \ill{} $(0 \leq i \leq n+1)$
\uncertain{gradués} de \uncertain{degré} \uncertain{un} les
$\partial_i^n$, \struck{\ill{}} avec les relations qu'on sait.
\struck{On} On \ill{} \ill{} y introduire \ill{} \ill{} les
\emph{relations}
\[
\partial_i^n = 0 \quad \text{pour} \quad 1 \leq i \leq n
\]
\marginal{$\partial_i^n : \Delta^n \to \Delta^{n+1}$, \quad $0 \leq i \leq n+1$}

de \uncertain{conditions} de \ill{} \ill{} $K^n$ \ill{} \ill{},
\ill{} \ill{} $K^i \to K^n$, \ill{} \struck{$K^i$} $i < n$ ;
$K^i \to K^n$ \ill{}
\[
\begin{array}{ll}
\partial_i^n K^{n-1} & 1 \leq i \leq n \\
u_2\, \partial_i^{n-1} K^{n-2} & 1 \leq i \leq n \\
\cdots & \\
u\ K^{0} &
\end{array}
\]
\note{ces lignes sont réunies par une accolade et entourées d'un ovale ;
un trait sinueux mène de la deuxième ligne vers l'intérieur de l'ovale, et
un trait oblique barre la première}

\ill{} opération \uncertain{simpliciale} \ill{}
\uncertain{correspondant} : \ill{} \ill{} \uncertain{croissante} \ill{}
$\Delta^{i}$ dans $\Delta^{n}$ \uncertain{qui} \ill{} \ill{}
\emph{\uncertain{sur} le \uncertain{dernier}}. Le \ill{} \ill{}
\[
K'^{n} \subset K^{n},
\]
et \ill{} \ill{} \ill{} \ill{} \ill{} \ill{} \ill{} \ill{} \ill{}
\uncertain{objet} \uncertain{simplicial} \ill{} $K^{*}$.
\note{la fin de la page, sous un trait, est une ligne en partie
soulignée dont seuls « objet simplicial » et « $K^{*}$ » se lisent}

\[
\text{\struck{$x_0 \leq x_1 \cdots \leq x_n$}} \qquad
(x_0, x_1, x_2) \longrightarrow (x_0, x_2)
\]
\note{la flèche descend ; au-dessus, un segment à points marqués}

\begin{tikzcd}[column sep=small, row sep=small, nodes={font=\scriptsize}]
 & & A_{[x_0 x_1 x_2]} & \\
 & & A_{[x_0, x_2]} \arrow[u] & \\
\Sigma A_{[x_0 x_1]} & \Sigma A_{[x_1, x_2]} & \Sigma A_{[x_0, x_1, x_2]} & A_{[x_0, x_2]} \arrow[l] \\
A_{[x_0]} \arrow[u] \arrow[ru] & A_{[x_1]} \arrow[u] & A_{[x_2]} \arrow[ru] &
\end{tikzcd}

\note{diagramme redessiné : sur la page les sommes sont écrites en ligne,
$\Sigma A_{[x_0 x_1]}\ \Sigma A_{[x_1,x_2]}\ \Sigma A_{[x_0,x_1,x_2]}$,
les $A_{[x_i]}$ en dessous avec de petites flèches montantes, dont une arche
courbe de $A_{[x_0]}$ vers la droite ; à côté, un croquis : trois droites
concourantes en un point, et deux points isolés ; les flèches exactes
entre $A_{[x_i]}$ et les sommes ne se laissent pas toutes attribuer}
\[
A_{[x_0 \ x_1}\ \cdots \qquad x_0, x_2 \qquad
A_{[x_0, x_2]} \qquad \partial^{1}_{1} A_{[x_0, x_2]}
\]
\[
\cdots \to K'' \to K' \to K \to 0
\]

\page{56}
\[
- \quad A_{[x_1 x_2, \ldots, x_n]} \longrightarrow
\prod_{\substack{x_0 \\ \text{précède} \\ x_1}} A_{[x_0, x_1, \ldots, x_n]}
\xrightarrow{\ d_0^{n}\ }
\prod_{\substack{x'_0, x_0 \\ (x'_0 \text{ précède } x_0)}} A_{[x'_0, x_0, x_1, \ldots, x_n]}
\]
\note{sous la première flèche, une flèche de retour marquée
$d_0^{n-1}$ ; « précède » est une lecture probable de l'abréviation ; dans
le premier indice, $x_1$ est surchargé ; à gauche, un segment à trois
points marqués}
\[
\left\lbrace
\begin{array}{l}
d'^{2} = 0 \\
d''^{2} = 0 \\
d'd'' = d''d'
\end{array}
\right.
\]
\note{ces relations sont entourées d'un cercle, sous un trait horizontal ;
à droite, un petit croquis}
\[
A_{[x_0, x_n]} \longrightarrow \text{\struck{$\prod$}}\ \prod A_{[x_0 x_i x_n]}
\]
\note{le premier produit est biffé avec son indice}
\[
\begin{array}{ccc}
A_{[x_0 x_{n-1}]} & & A_{[x_1, x_n]} \\
\swarrow \ \searrow & & \swarrow \ \searrow \\
A_{[x_0, x_{n-1}, x_n]} & & A_{[x_0, x_1, x_n]}
\end{array}
\]
\note{à droite de ce schéma, séparé par un trait vertical : « $x_0$
\uncertain{min} \,|\, $x_n$ \uncertain{maximal} », et dessous
« \uncertain{non} \ill{} » ; un trait courbe sépare le tout du bas de
la page}
\[
\begin{array}{llll}
X & \prod \hat{\underline{\mathcal{O}}}_{x} & & \prod G(\hat{\underline{\mathcal{O}}}_{x}) \longleftarrow \cdot \\
C & \prod_{\mathrm{loc}} \hat{K}_{x} & K & \prod_{\mathrm{loc}} G(\hat{K}_{x}) \longleftarrow G(K)
\end{array}
\]
\note{de $\prod G(\hat{\underline{\mathcal{O}}}_x)$ et du point à sa droite
descendent deux petites flèches vers la seconde ligne}
\[
\begin{array}{c}
G(\hat{A}) \\
\downarrow \\
G(\hat{K}) \longleftarrow G(K)
\end{array}
\qquad\qquad H
\]
\[
K \to K/A, \qquad G(K) \qquad G(K)/G(A)
\]

\page{57}
\note{en haut, un treillis des $A_{[\ldots]}$ pour la chaîne
$x_0 < x_1 < x_2 < x_3$, dessiné en étages avec de petites flèches
montantes ; nous en donnons les étages, de haut en bas ; trois longues
flèches au crayon rouge renvoient des termes de droite vers ceux de
gauche}
\[
\begin{array}{l}
A_{[x_0 x_1 x_2 x_3]} \\
A_{[x_0 x_1 x_2]} \quad A_{[x_1 x_2 x_3]} \quad \text{\struck{$A_{[x_1 x_2 x_3]}$}} \quad A_{[x_0 x_2 x_3]} \quad A_{[x_0 x_1 x_3]} \\
A_{[x_0 x_1]} \quad A_{[x_1 x_2]} \quad A_{[x_2 x_3]} \quad \text{\struck{$A_{[x_2 x_3]}$}} \quad A_{[x_0 x_2]} \quad A_{[x_0 x_1]} \quad \text{\struck{$A_{[x_2 x_3]}$}} \quad A_{[x_1 x_3]} \quad \text{\struck{$A_{[x_1 x_3]}$}} \quad \text{\struck{$A_{[x_2 x_3]}$}} \\
A_{[x_0]} \quad A_{[x_1]} \quad A_{[x_2]} \quad A_{[x_3]}
\end{array}
\]
\note{les termes biffés le sont sur la page ; plusieurs indices sont
surchargés et leur lecture reste incertaine}
\[
\overline{A_{[x_0 x_2 x_3]}} \cap \overline{A_{[x_0 x_1 x_3]}}
\]
\[
\overline{\overline{A_{[x_0, x_3]}}}
\]
\note{la seconde ligne est écrite sous la première, sans signe entre
elles}
\[
0 \qquad (A_{[x_0 x_2 x_3]} + A_{[x_0 x_1 x_3]})
\]
\[
A_{[x_0 x_2]} \times A_{[x_1, x_3]} \times \bigl(A_{[x_0 x_2 x_3]} \times A_{[x_0 x_1 x_3]}\bigr)
\]
\[
0 \qquad A_{[x_0 x_2]} \qquad A_{[x_{\text{\uncertain{$2$}}} x_3]} \qquad A_{[x_{\text{\uncertain{$0$}}} x_3]}
\]
\note{des flèches croisées relient chaque terme de la dernière ligne à
deux facteurs du produit ; le produit entre parenthèses est entouré, et une
flèche en part vers la formule :}
\[
\mathrm{Im}\, A_{[x_0 x_3]} + \mathrm{Im}\, A_{[x_1, x_3]}
\qquad
\alpha \quad \beta \quad -\bar{\alpha} \quad -\bar{\beta},
\qquad \alpha \circ \beta
\]
\[
\begin{array}{llllll}
S_0 & S_1 & S_2 & S_{n+1} \xleftarrow{\ d_i^{n}\ } S_n & & \\
 & & & d_i^{n}\sigma \quad \ \sigma & & \\
 & & & A_{d_i^{n}\sigma} \quad A_\sigma & &
\end{array}
\]
\note{sous $S_n$, le signe $\in$ couché ; un trait horizontal traverse
ensuite la page}
\[
K \qquad d_i^{n} \quad 0 \leq i \leq n+1 \qquad\qquad d' \pm d''
\]
\[
\bigcap_{1 \leq i \leq n+1} \mathrm{Ker}\, d_i^{n}
\qquad
\text{\struck{$\ill{}$}}\ \mathrm{Coker}\, d_i^{n} \quad \underline{1 \leq i \leq n}
\]
\[
d_i^{n+1} d_0^{n} x = 0 \quad \text{si} \quad 1 \leq i \leq n+2
\]
\note{à droite, des croquis de segments subdivisés, légendés $\Delta^n$,
$\Delta^{n+1}$, $\Delta^{n+2}$ et « $x_0\ x_1\ x_2 - x_n$ »}
\[
A_{[x_0, x_2]} \xleftarrow{\ d\ } A_{[x_0, x_1, x_2]}
\]
\note{le « $d$ » est écrit sous la flèche}

1°) $\bigcap \mathrm{Ker}$ \emph{des \uncertain{dégénérescences}}
\note{la ligne s'arrête là, au bas de la page}

\page{58}
\note{feuillet écrit dans le sens de la largeur ; un quadrillage au crayon
(trois colonnes, trois lignes, étiquetées d'abréviations illisibles, en
haut « \ill{} » et « \ill{} », à gauche « \ill{} ») porte le diagramme
ci-dessous ; une diagonale au crayon le traverse ; dans les cases vides,
des « $0$ » ; nous notons $\mathfrak{r}$ l'idéal premier écrit d'une lettre
gothique que nous ne lisons pas sûrement}

\begin{tikzcd}[column sep=small, row sep=small, nodes={font=\scriptsize}]
\hat{A} \arrow[d] & & \\
\varprojlim \hat{A}/\mathfrak{r}^{n}\hat{A} \otimes_{A} A_{\mathfrak{r}}/\mathfrak{r}^{n} A_{\mathfrak{r}} \arrow[d] & \widehat{A_{\mathfrak{r}}} \arrow[l] \arrow[d] & \\
\prod_{\mathfrak{r}} A_{[\mathfrak{m}, \mathfrak{r}, 0]} & \hat{A}_{\mathfrak{r}} \otimes_{A} K = \widehat{K_{\mathfrak{r}}} \arrow[l] & K \arrow[l]
\end{tikzcd}

\note{sous le terme de gauche, une accolade :
$\hat{A} \otimes_{A} A_{\mathfrak{r}}/\mathfrak{r}^{n} \hat{A}_{\mathfrak{r}}$,
puis « $A_{[\mathfrak{m},\mathfrak{r}]}$ » ; sous
$\widehat{K_{\mathfrak{r}}}$, « $A_{[\mathfrak{r}, 0]}$ » ; sous le
quadrillage : « \ill{} $\supset \mathfrak{r} \supset 0$ »}
\[
\mathrm{Hom}(M, A), \quad \mathrm{Ext}^{1}(M, A) \cdots
\qquad\qquad
\mathrm{Hom}(K_{\mathfrak{p}}, K_{\mathfrak{q}}) = A_{\mathfrak{q},\mathfrak{p}} \quad (\mathfrak{q} \supset \mathfrak{p})
\]
\note{la première formule est soulignée d'une accolade}
\[
\begin{array}{l}
S_{A}(M) \\
\mathrm{Hom}_{A}(M, \mathcal{A}) \\
\mathrm{Ext}^{*}(M, A) \\
\Uparrow \\
H^{*}\, \mathrm{Ext}^{*}(M, \mathcal{A}) \\
\mathrm{Ext}(\hat{M}, A) \qquad \hat{A}\ \hat{K}\ K
\end{array}
\]
\note{en bas, un second quadrillage au crayon porte le même diagramme
pour les groupes multiplicatifs :}

\begin{tikzcd}[column sep=small, row sep=small, nodes={font=\scriptsize}]
\hat{A}^{*} \arrow[d] & & \\
\varprojlim (\hat{A} \otimes_{A} A_{\mathfrak{r}}/\mathfrak{r}^{n} A_{\mathfrak{r}})^{*} \arrow[d] & \widehat{A_{\mathfrak{r}}}^{*} \arrow[l] \arrow[d] & \\
\cdot & \widehat{K_{\mathfrak{r}}}^{*} \arrow[l] & K^{*} \arrow[l]
\end{tikzcd}

\[
A \ \ill{} \ill{}
\qquad
0 \to A \to \mathcal{A} \to \mathcal{A}_1 \to 0
\]
\[
\mathrm{Ext}^{1}(M, A) \to \mathrm{Ext}^{\text{\uncertain{$1$}}}(M, \mathcal{A}) \to \ill{}
\]
\note{la dernière formule est doublement soulignée ; à droite, dans un
cadre, des $\mathfrak{p}$ surchargés et une ligne raturée}

\page{59}
\note{la page commence au milieu d'une phrase ; elle semble continuer la
page 60, dont la dernière ligne annonce des accouplements
$A_{c'} \times A_c \to A_{c'c}$ (les deux feuillets seraient intervertis)}
[\uncertain{dans} \ill{} \ill{} \ill{} \ill{} \ill{} $(c_i)$ \ill{}
\ill{} \ill{} $x_0 \geq \cdots \geq x_n$ \ill{} \ill{} \ill{} $(c_i)$,
\ill{} \ill{} \uncertain{chaîne} \uncertain{saturée} \ill{}
$x_0 \geq \cdots \geq x_n$ \ill{} $\in A_c$]
\note{interligne serré, surchargé et en partie biffé ; au-dessus,
« $X$ est une \uncertain{catégorie} \ill{} » et « $A_c$ » biffé}
une \emph{multiplication}. L'élément unité \ill{} \ill{} \ill{} \ill{}
\uncertain{de} $\struck{A_c}$ \ill{} \ill{} \ill{} \ill{}
\uncertain{compatibles}, la \uncertain{somme} des $A_c$ devient un
\uncertain{anneau} $A^{*}$ [\ill{} \uncertain{gradué} \ill{}
\emph{\uncertain{gradué}} par la \uncertain{longueur} des
\uncertain{chaînes}, \uncertain{non} \uncertain{commutatif}],
\struck{\ill{}} la \uncertain{somme} des $A_x$ devient un
\uncertain{sous-anneau} (\uncertain{commutatif}), \uncertain{noté}
$A^{0}$. $A^{1}$ est la somme des $A_{[x_0, x_1]}$, \uncertain{soit}
$\Theta$ l'élément de $A^{1}$ dont les composantes sont
$1_{[x_0, x_1]}$. \struck{Combien est $\Theta^2$ ?} On a
\[
\Theta^{2} = \sum_{\substack{(x_0, x_1) \\ (y_0, y_1) \\ \text{chaînes saturées}}} 1_{[x_0, x_1]} \cdot 1_{[y_0, y_1]}
= \sum_{\substack{[x_0, x_1, x_2] \\ \text{chaînes saturées}}} 1_{[x_0, x_1, x_2]}
\]
On va maintenant diviser par l'idéal \uncertain{bilatère}
\note{« bilatère » est ajouté au-dessus de la ligne}
\uncertain{engendré} par $\Theta^{2}$,
\struck{[\ill{} \ill{} \ill{} \ill{} \ill{}]}
\struck{\ill{} \ill{}}, \ill{} \ill{} \ill{} \ill{} \ill{}.
En d'autres termes, \ill{} \ill{} \ill{} \ill{}, \ill{}
$u_1 \pm u_2$ \ill{} \ill{} \ill{} il \ill{} \uncertain{négliger} les
\[
\sum_{x_1} \varphi \qquad \Bigl(\ill{}\ \delta(a, b) = 2\Bigr)
\]
\note{sous la somme, deux lignes biffées et illisibles}
\note{à droite, un croquis : un segment de $a$ à $b$, coupé entre $\alpha$
et $\beta$ par un trait oblique légendé « \uncertain{fenêtre} » ;
dessous :}
\[
\begin{array}{ccc}
A_{[a - \alpha, \delta, \beta, \ldots, b]} & A_{[a - \alpha]} & A_{[\beta, \ldots, b]} \\
\uparrow \varphi_{\sigma} & \uparrow & \uparrow \\
A_{[\alpha - a]} \otimes A_{[\beta \ldots \beta]} & A_{[\alpha]} & A_{[\beta]}
\end{array}
\qquad
[A_{\alpha \ \ill{} \ b}]
\qquad
x_0 \longrightarrow x_n
\]
\note{sous le $\otimes$, « $A_{\ill{}}$ » et « $b - \beta$ » ; les indices
de ce groupe sont d'une lecture très incertaine}

\page{60}
$I$ ensemble ordonné satisfaisant à la « condition des chaînes »
[a) La longueur des chaînes joignant $x \geq y$ est bornée
b) Deux chaînes saturées joignant $x$ à $y$ ont même longueur
$d(x, y)$].
\note{les conditions a) et b) sont écrites en interligne, au-dessus de la
ligne suivante ; le « $I$ » initial est surchargé}

Pour toute chaîne \emph{saturée}
\[
c = (x_0 \geq x_1 \geq \cdots \geq x_n),
\]
on \uncertain{suppose} donné un \emph{anneau}
\[
A_c = A_{[x_0, \ldots, x_n]}
\]
et pour une chaîne de longueur \uncertain{maximum} (\uncertain{supposée}
$n \geq 1$) des \emph{hom.\ d'anneaux}
\[
\begin{array}{ll}
u_1 : A_{[x_0, \ldots, x_{n-1}]} \longrightarrow A_{[x_0, \ldots, x_n]} & [\text{i.e.\ } (1_{x_n}\Theta)_0] \\
u_2 : A_{[x_1, \ldots, x_n]} \longrightarrow A_{[x_0, \ldots, x_n]} & [\text{i.e.\ } d(\Theta 1_{x_0})]
\end{array}
\]
\note{les gloses entre crochets, à droite, sont d'une autre encre ;
leur lecture est incertaine}
[on note que
\[
u_1, u_2 : A_{c'} \longrightarrow \prod_{\substack{c \supset c' \\ \lg c = \lg c' + 1}} A_c
\]
]
\note{suit une ligne biffée et illisible}

on suppose que $u_1$ et $u_2$ commutent : si $n \geq 2$

\begin{tikzcd}
A_{[x_1, \ldots, x_{n-1}]} \arrow[r, "u_2"] \arrow[d, "u_1"'] & A_{[x_0, x_1, \ldots, x_{n-1}]} \arrow[d, "u_1"] \\
A_{[x_1, \ldots, x_n]} \arrow[r] & A_{[x_0, x_1, \ldots, x_{n-1}, x_n]}
\end{tikzcd}

de sorte que cela permet de définir sans \uncertain{ambiguïté} des
\uncertain{hom.}\ d'anneaux
\[
u_{c,c'} : A_{c'} \longrightarrow A_{c} \quad \text{si} \quad c \supset c'
\]
[\struck{\ill{}} et \ill{} \ill{} \ill{} une condition de
\emph{transitivité}, qui exprime que $(A_c)$ est un système
\uncertain{inductif} sur l'\uncertain{ens.}\ des chaînes ordonnées. De
\uncertain{plus}, \ill{} \ill{} \ill{} \emph{\uncertain{multiplication}}
dans les $A_c$, on se \uncertain{donne} des
\struck{\ill{}} \uncertain{accouplements} \uncertain{bi-additifs}
\[
A_{c'} \times A_{c} \longrightarrow A_{c'c}
\]
si $c', c$ sont \uncertain{compatibles} dans un certain sens, et
\uncertain{satisfont} les
\note{la page s'arrête sur ce mot ; la suite semble être en tête de la
page 59}

\end{document}
